\subsection{The complete three-pointed torus word dictionary}

\begin{result}{D.29}{Proved}{All 108 cells and 304 face maps for $(1,3)$}
The literal $P_2(1,3)$ complex of groups is given by
\path{code/d_p2_elliptic_three.json}; its complete readable
version is \path{code/d_p2_elliptic_three_dictionary.md}.
There are $22$ vertex orbits, $53$ edge orbits and $33$
triangle orbits.  The barycentric diagram has $108$ objects,
$304$ face homomorphisms and $198$ coherence triangles.
Every local group has a named presentation; every generator
image and every twisting element is a word in its named
target generators.  In total the executed certificate checks
$922$ generator images, $126$ defining relators in $28$
local presentations, and $594$ generator equations for the
$198$ coherences.  Exactly seven coherence triangles have
a nonidentity twist.

With the specified $107$-edge maximal tree, its fundamental
presentation is explicitly isomorphic to (D.22).
The inverse sends $a_1,a_2,a_3$ to the corresponding node
generators at the three-curve vertex, and $b$ to the named
$b$ at the all-three-points separating vertex.
The open stratum supplies all $198$ cone tetrahedra,
with the two sides recorded as exact words.  This gives
the fifth requested group-presentation dictionary.
The further compact Peiffer comparisons are supplied
with their complete cellular maps in D.38.
\end{result}
\begin{proof}
We specify the remaining stabilizers, all choices in the
finite diagram, and its inverse presentation map.
Throughout the letters are the actual twists of (D.22).
Put
\[
 h=(a_1b)^3,\quad d_0=h^2,\quad
 \widetilde d_1=d_0d_3^{-1}d_2^{-1},\quad
 \widetilde d_2=d_2,\quad\widetilde d_3=d_3.
                                                               \tag{D.24}
\]
The regular neighborhood of $a_1\cup b$ is a one-holed
torus.  Its complementary four-holed sphere contains
all three marked points.  Soroko's Figure~2 identifies
$d_2,d_3$ with the two standard intersecting pair curves
in that sphere; their intersection number is two.
The boundary-framed lantern relation makes
$\widetilde d_1$ its third pair twist.  The chosen order
in (D.24) is certified by the exact identities
\[
 q_2d_1q_2^{-1}=\widetilde d_1,\qquad
 q_3d_1q_3^{-1}=d_0d_2^{-1}d_3^{-1}.                  \tag{D.25}
\]
Both were checked independently by
\path{code/b_d4_tail.py} using D.28's full word algorithm.
Thus no choice of an unrecorded third pair curve is needed.

Here is the full list of chosen multicurve representatives.
For each $i$, write $\{j,k\}=\{1,2,3\}\setminus\{i\}$
with $j<k$.  In the $L_i$ row use instead the distinguished
ordered pairs
\[
             (j_i,k_i)=(3,2),(1,3),(1,2)
                       \quad(i=1,2,3).
\]
\[
\begin{array}{c|c|c}
\text{name}&\text{multicurve}&\text{number of orbits}\\\hline
N&\{a_1\}&1\\
S_0&\{d_0\}&1\\
S_i&\{d_i\}&3\\
L_0&\{a_1,d_0\}&1\\
L_i&\{a_{j_i},d_i\}&3\\
C_i&\{a_j,a_k\}&3\\
T_i&\{d_0,\widetilde d_i\}&3\\
F_i&\{a_1,d_0,\widetilde d_i\}&3\\
E_i&\{a_j,a_k,d_i\}&3\\
\Theta&\{a_1,a_2,a_3\}&1
\end{array}                                                     \tag{D.26}
\]
Their dual graphs, with every puncture label retained, are
exactly the $22$ nonempty stable graphs of D.28.
The JSON specifies the ordered edge list for every row,
including which edge is the loop and which is the pair
separator.  These orderings determine the masks used below.

\emph{The nonseparating stabilizer.}
Cut $a_1$ and label the four finite points of the resulting
sphere as
\[
         (1,2,3,4)=(\text{node}_1,\text{node}_2,p_3,p_1),
                    \qquad\infty=p_2.
\]
Let $B_4(\mathfrak S_2)$ be the subgroup of $B_4$ whose
permutations belong to $\{1,(12)\}$, with standard pure
generators $A_{ij}$ and the half braid $s=\sigma_1$.
Its presentation is the pure four-braid presentation,
$s^2=A_{12}$, and the six actual conjugation formulas
$sA_{ij}s^{-1}=\sigma_1A_{ij}\sigma_1^{-1}$.
Those right sides are individually expanded in the
certificate, by the exact braid combing of B.18.
Write
$Z_4=A_{12}A_{13}A_{23}A_{14}A_{24}A_{34}$.
Then
\[
 G_N=\bigl(B_4(\mathfrak S_2)\times\mathbb Z(a_1)\bigr)
                   /\langle Z_4a_1^2=1\rangle.          \tag{D.27}
\]
Its literal generator embedding is
\[
\begin{aligned}
 s&=h a_1^{-1},& A_{12}&=s^2,\\
 A_{13}&=a_2a_1^{-1},& A_{14}&=a_2^{-1}a_3d_2^{-1},\\
 A_{34}&=d_2,& A_{23}&=sA_{13}s^{-1},\qquad
 A_{24}=sA_{14}s^{-1}.
\end{aligned}                                                     \tag{D.28}
\]
To prove this identification, use zero-rotation
parameterizations of the two node holes.  A pure braid
is lifted by rigid translation of disjoint sufficiently
small discs; their radii can be chosen continuously from
the minimum point separation.  This gives the specified
zero framing.  Sewing identifies both node-hole boundary
twists with $a_1$.  Capping the outer boundary gives
$Z_4(z_1z_2)=1$, hence precisely (D.27).
The half exchange has image $h a_1^{-1}$; its square is
$d_0a_1^{-2}$.  A twist enclosing node $1$ and $p_3$ has
image $a_2a_1^{-1}$, and the pair $p_3,p_1$ has image
$d_2$.  The three-point cluster $\{1,3,4\}$ has full
twist $A_{13}A_{14}A_{34}$ and image $a_3a_1^{-1}$;
solving yields the stated $A_{14}$.  Conjugation by the
actual half exchange gives the other two generators.
These generate the entire cut-sphere group, and the
only sewing kernel is the displayed boundary relation.
This proves both surjectivity and injectivity of (D.27),
including the branch-exchanging coset.
Useful inverse words, also checked exactly, are
\[
\begin{aligned}
 a_2&=A_{13}a_1,&a_3&=a_2A_{14}A_{34},&h&=sa_1,\\
 d_3&=a_1^2A_{12}A_{13}A_{23},&
 \widetilde d_1&=a_1^2A_{12}A_{14}A_{24}.&&
\end{aligned}                                                     \tag{D.29}
\]

\emph{The remaining vertex groups.}
The four families $S_i,C_i,E_i,\Theta$ are proved in
D.28.  The other groups have the following full
presentations, with literal generator names:
\[
\begin{aligned}
 G_{S_0}&=B_3(a_1,b)\times F(d_2,d_3),\\
 G_{L_0}&=\mathbb Z^2(a_1,h)\times F(d_2,d_3),\\
 G_{T_i}&=B_3(a_1,b)\times\mathbb Z(\widetilde d_i),\\
 G_{F_i}&=\mathbb Z^3(a_1,h,\widetilde d_i),\\
 G_{L_i}&=\mathbb Z(a_{j_i})\times
       \langle a_{k_i},r_i\mid[r_i^2,a_{k_i}]=1\rangle,\\
 r_i&=a_{k_i}a_{j_i}^{-1}(a_{j_i}b)^3,
                       \qquad r_i^2=d_i.
\end{aligned}                                                     \tag{D.30}
\]
For $S_0$, cut the one-holed torus from the three-puncture
disc.  The latter's pure braid group has central boundary
$\mathbb Z(d_0)$ and free complementary factor
$F(d_2,d_3)$.  Sewing identifies $d_0$ with $h^2$ in
$B_3$ and eliminates this one central generator; it gives
exactly the first product.  Stabilizing $a_1$ in its
$B_3$ factor gives $\mathbb Z^2(a_1,h)$, as in B.10.
Stabilizing a specified pair curve in the four-holed
sphere replaces its free group by that curve's cyclic
stabilizer.  This proves the next two products and the
$F_i$ row, with the node meridian $h^2$ retained.
Finally $L_i$ is the boundary-retaining nonseparating
stabilizer inside $B_4(a_{j_i},b,a_{k_i})$.
B.10's cutting proof gives
$h_i=a_{j_i}a_{k_i}^{-1}r_i$ and $r_i^2=d_i$;
it gives the last presentation without imposing a
spurious order-two relation on $r_i$.
The program checks each relation and each displayed
embedding in the exact group (D.22).

\emph{Flag stabilizers and all transports.}
The only nontrivial permutations of distinct edges of
these representatives occur at $C_i,E_i$: $q_i$ exchanges
the two nonseparating edges.  Their ordered subgroups are
\[
 G_{C_i}^{\mathrm{ord}}
    =\mathbb Z^2(a_j,a_k)\times F(a_i,d_i),\qquad
 G_{E_i}^{\mathrm{ord}}=\mathbb Z^3(a_j,a_k,d_i).          \tag{D.31}
\]
A flag has the full group of its largest multicurve
unless one of its subsets distinguishes these two edges;
in that case it has the appropriate ordered subgroup.
Loop-branch reversals in $N,L_i,L_0,F_i$ remain in their
already displayed groups; they do not permute distinct
edges.  Thus (D.27), (D.30), D.28 and (D.31) give all
$28$ presentations used by the $108$ objects.

Here is the completely deterministic representative
rule used in the supplied tables.  For each largest
multicurve, take every decreasing chain of nonempty
edge masks and choose the lexicographically least one
in its actual edge-permutation orbit.  This gives the
$22,53,33$ orbits.  For a face, first take its largest
retained mask.  Among
\[
             1,q_1,q_2,q_3,q_1^{-1},q_2^{-1},q_3^{-1}
\]
choose the first word carrying its individual twist
set to the chosen representative in (D.26).
Such a word exists for every subset in the executed
list.  Transport the remaining masks by its actual
edge bijection.  If they are not the chosen least
flag, multiply on the left by the target $q_i$ that
puts them there.  This is possible precisely in the
$C_i,E_i$ cases.  The resulting word is $u_a$ for the
face arrow $a$.

These are actual transporters.  The identities are
checked by conjugating each twist and using D.28's
full word problem.  Equality of two positive Dehn
twists identifies their curves: if $T_c=T_d$, applying
them to $c$ and using
$i(c,T_d(c))=i(c,d)^2$ first makes the curves disjoint;
independence of the twists of a multicurve then makes
them equal.  The ordered masks consequently identify
the entire face flag, not only its abstract dual graph.
The face injection is exactly
$\psi_a(g)=u_a g u_a^{-1}$.

Every such injection is rewritten in the target's
named generators.  Equations (D.28)--(D.30) provide the
needed inverse substitutions at $N$ and $L_i$.
Conjugations of the remaining leaves are obtained from
\[
 q_i a_i q_i^{-1}=a_ja_k a_i^{-1}d_i^{-1},\qquad
 q_i^{-1}a_iq_i=a_ja_k d_i^{-1}a_i^{-1},
\]
while $q_i$ exchanges $a_j,a_k$ and fixes $b$.
Together with (D.25), these formulas cover every
transport in the list.  Each of the $922$ resulting
words is independently evaluated in (D.22) and
checked against $u_a g u_a^{-1}$.  The faithful
geometric embeddings of the local groups then prove
that these are homomorphisms into the named target
presentations, with no kernel introduced by rewriting.

For composable face arrows $a,b$, the twisting word is
\[
                    f_{b,a}=u_bu_a u_{ba}^{-1},\qquad
             \psi_b\psi_a=\operatorname{Ad}(f_{b,a})\psi_{ba}.
\]
For clarity the nonidentity
values can be listed completely.  A row specifies the
triangle flag and the intermediate edge flag; the
last column gives its target-group element.  The
subsets below are the actual curves in
the largest representative before its recorded face
transport.
\[
\begin{array}{c|c|c|c}
\text{triangle}&\text{intermediate edge}&\text{target}&f_{b,a}\\\hline
(F_1,\{a_1,\widetilde d_1\},\widetilde d_1)
 &(L_1,d_1)&S_1&a_3^{-1}b^{-1}a_3^{-1}a_2ba_3\\
(E_1,\{a_2,d_1\},a_2)&(L_1,a_3)&N&h d_3^{-1}\\
(E_1,\{a_2,d_1\},d_1)&(L_1,d_1)&S_1&q_1\\
(\Theta,C_1,a_3)&(C_1,a_2)&C_1&q_1\\
(\Theta,C_1,a_3)&(C_1,a_2)&N&h d_2^{-1}\\
(\Theta,C_2,a_3)&(C_2,a_1)&C_2&q_2\\
(\Theta,C_3,a_2)&(C_3,a_1)&C_3&q_3
\end{array}                                                     \tag{D.32}
\]
The fourth and fifth rows use the two possible final
vertices of the same intermediate edge.  The words
at $N$ are read by (D.29), and $q_1$ at $S_1$ is its
ordinary four-braid half twist.  Their exact global
checks are
\[
\begin{aligned}
 q_3q_1q_2^{-1}&=h d_2^{-1},&
 q_2q_1q_3^{-1}&=h d_3^{-1},\\
 q_2^{-1}q_3^{-1}&=a_3^{-1}b^{-1}a_3^{-1}a_2ba_3.&&
\end{aligned}
\]
All other $191$ twisting elements are identity.
The $198$ equations are checked on all $594$ source
local generators in the certificate.  This establishes
the entire complex of groups with its actual twists.

\emph{The presentation maps and dimension-three words.}
Select the first spanning tree in the given arrow order
using only $u_a=1$ arrows; the executed list has $107$
edges and connects all $108$ objects.  The resulting
finite presentation has the object-group generators
and arrow letters $t_a$, relations
\[
 t_a g t_a^{-1}=\psi_a(g),\qquad
 t_b t_a=f_{b,a}t_{ba},\qquad
 t_a=1\ \text{on the chosen tree}.
\]
The map sends $t_a$ to $u_a$; together with the literal
local embeddings this defines its map to (D.22).  The diagram is the
actual action complex of groups on $P_2(1,3)$: its
cell representatives and stabilizers were identified
above, including every parallel-edge incidence.
D.6 proves that this development is simply connected.
Lifting edge paths and elementary face homotopies
therefore proves that the displayed fundamental-group
map is an isomorphism.  Its inverse on the four
presented global generators is the one stated in the
result, since their local images are literally
$a_1,a_2,a_3,b$.  The isomorphism proves that these
specified inverse words satisfy all six Artin
relations and the central relation of (D.22).

Finally adjoining the empty multicurve is coning the
whole development.  Set the transporter from every
object to this open vertex equal to $1$.  The twist
on an object--face--open triangle is then precisely
its already specified $u_a$.  For each of the $198$
three-step flags the tetrahedral identity becomes
\[
                       u_bu_a=f_{b,a}u_{ba}.
\]
Both sides agree as freely reduced words, not only as
elements of the presented group.  These are all the
cone tetrahedra for the unique dimension-three
stratum of $M_{1,3}$.  Their further expansion in a
chosen compact presentation's free crossed module
is completed with its exact cellular comparison in D.38.
\end{proof}
